\section{Sampling and Regression to Reduce Experiment Time}
\label{sec:regression}

\vspace{-.1in}
\paragraph{Cost of extensive evaluation.}
Given all the settings to measure and the length of each experiment,
one can estimate the total machine hours necessary to run all experiments
once. In our experience, two factors are adding uncertainties to such
estimation: first, both TPC-C and YCSB have an initialization phase to
populate data, which could take long, before running any experiments. 
%For example, using a single client to populate a large
%amount of data may cause the initialization phase to be longer than the
%following experiments. 
Second, it is generally a recommended practice to
run experiments multiple times to check whether the result is stable. In
our experience, we find most results are stable so we decide to run three
times, but we do observe a few systems have unstable performance under
a subset of settings. In general, we observe unstable performance
is often an indicator of underlying issues, so we should investigate the
reason instead of simply running experiments more times.

Table~\ref{tab:summary} summarizes the settings we finish to run
and the total machine hours we use to run each experiment once. In general, we observe that,
perhaps not surprisingly, distributed systems require much more
machine hours since they not only add a parameter
(i.e., number of servers), but also require more machines to run these distributed experiments.
As a result, for single-node systems, we find the total machine hours is well affordable,
but for distributed systems, we have to carefully pick parameter values to minimize
total machine hours. To be concrete, 1800 machine hours means 75 machines
for a day; considering we need to run experiments three times and
may need to run multiple iterations due to optimization or bug fix, we may need 75
machines for a week or longer, which is still manageable but could be a burden.
If one needs to tune more parameters or more values of
a parameter, the total machine hours required can quickly get out of control. 

% Miao: This can enable line break within a cell
% gives us better control of when to start new line
\newcommand{\specialcell}[2][c]{%
  \begin{tabular}[#1]{@{}c@{}}#2\end{tabular}}
 
% \specialcell{Foo\\bar}     vertically centered
% \specialcell[t]{Foo\\bar}  aligned with top rule
% \specialcell[b]{Foo\\bar}  aligned with bottom rule



\begin{table*}[ht]
    \centering
    \begin{tabular}{c|c|c} System    &  Settings & Machine hours  \\
    
    % Miao: Competitors run time are included
     Calvin-TPCC     &  \#WH=\{1,2,4\dots128\},\#Node=\{1,2,4\dots32\},Cross-WH(\%)=\{0,10\dots100\}   &  832 \\
     \cline{1-3}
     Silo-TPCC       & \specialcell{\#WH=$\{1,2,4 \dots 512\}$, \#Workers=\{1,2,4,8,16,28\},\\     \#Txn-Types=\{5,2,NewOrder-only\},Cross-WH-NewOrder(\%)=\{0,10 \dots 100\}}  & 130 \\
     \cline{1-3}
     DrTM-TPCC       &   \specialcell{\#Node=\{2,4,8\}, \#Workers=\{2,4,8,16,32\}, \\ 
     \#WHperNode=$\{1,2,4 \dots 256\}$, Cross-WH(\%)=\{0,20 \dots 100\}}  & 192                           \\
     \cline{1-3}
     Janus-TPCC      & \specialcell{ \#Clients=\{1,10,20,40,80,100,200,500,1000,2000,4000\}\\ \#Replica=\{1,3,5,7\}, \#CPUperNode=\{1,2,4,8,16\},\\ \#WH=$\{1,2,4\dots\mid \#WH\times\#Replica/\#CPUperNode \leq 25\}$ } & 1886 \\
     \cline{1-3}
     Cicada-TPCC     &    \#WH=\{1,2,4\dots256\},\#Workers=\{1,2,4,8,16,28\},\#Txn-Types=\{5,2\}  & 9 \\
     \cline{1-3}
     GAM-TPCC        & \specialcell{\#WH=\{2,4,8,16,32 / Node\}, \#Node=\{2,4,8\}, \\ Cross-WH(\%)=\{0,20,40\dots100\},
     \#Workers=\{1,2,4,8,16\},\#Txn-Types=\{2\}   } &     721                   \\
     \cline{1-3}
     Star-TPCC      &      \specialcell{\#WH=\{16,32,128,256\},\#Workers=\{1,2,4,8,16\},\\Cross-WH(\%)=\{0,20,40\dots100\},\#Node=\{2,4,8\}, \#Txn-Types=\{2\}  }   & 210                  \\
     \cline{1-3}
     Aria-TPCC      &  \specialcell{ \#WH=\{1,2,4\dots1024\},Cross-WH(\%)=\{0,10\dots100\},\#Nodes=\{1,2,4,8\}  \\ \#Workers=\{1,2,4,8,12\},\#Locks=\{1,2,4,6|\#Locks<\#Workers\}}   & 360\\
     \cline{1-3}
     Star-YCSB      &   \specialcell{\#KV Pairs=\{3.2M, 6.4M, 25.6M, 51.2M \}, \#Workers=\{2,8,16\}, \\Cross-partition(\%)=\{0,20,40$\dots$100\},  Zipf$\theta$=\{0,0.4,0.8,0.99\}, \\R/W Ratio=\{0/100,20/80,40/60,100/0\}, \#Node=\{2,4,8\} }   & 384                     \\
     \cline{1-3}
     Silo-YCSB      & \specialcell{\#KV Pairs=\{128K,512K,2M, 8M, 32M,128M\} \\ \#Workers=\{1,2,4,8,16,28\},R/W Ratio=\{0/100,20/80,40/60,100/0\}} & 5\\
     \cline{1-3}
     Cicada-YCSB    &    \specialcell{\#KV Pairs=\{128K,512K,2M, 8M, 32M\},KV size=\{8B,32B,128B, 1KB\},\\ \#Workers=\{1,4,16,28\}, Zipf$\theta$=\{0,0.4,0.8,0.99\},\\ R/W Ratio=\{0/100,20/80,40/60,100/0\},batch size=\{1,10,100\}} & 300 \\
     \cline{1-3}
     HERD-YCSB      &    \specialcell{\#KV Pairs=\{6M, 24.5M, 96M, 384M, 1536M\}, KV size=\{8B, 32B, 128B, 512B, 1KB\},\\ Zipf$\theta$=\{0,0.4,0.8,0.99\}, R/W Ratio=\{0/100,20/80,40/60,100/0\}} & 325                       \\
     \cline{1-3}
     TAPIR-YCSB     &   \specialcell{\#Clients=\{5,10,15,20,25\},\#KV Pairs=\{128K,512K,2M, 8M, 32M\},\\R/W Ratio=\{0/100,20/80,40/60,100/0\},\\ (\#shard,\#replica)=\{(1,3),(1,5),(1,7),(2,3),(2,5),(3,3)\} } & 1218      \\
     
    \end{tabular}
    \caption{The settings and machine hours we use for each system, including its competitors. Some systems don't support the tuning of a certain parameter. We set the length of each experiment to
    be the same as the one used by the original paper.}
    \vspace{-.1in}
    \label{tab:summary}
\end{table*}

\vspace{-.1in}
\paragraph{Sampling and regression.}
We study whether we can reduce the total machine hours with the classic
sampling and regression idea: if we could run a subset of experiments and use the results
to accurately predict the performance under the remaining settings, then we don't need to
run all experiments.

We use the following methodology: for each system, from the results we have collected, we choose a subset of them as the training
set and use the remaining results as the testing set; we train a model on the training
set and test its accuracy on the testing set; we further tune the size of the training set to see how
it affects the accuracy of prediction. 
We used scikit-learn~\cite{scikit-learn} for our data analysis.
%\todo{Yujie: When choosing the training set, we have a constraint that the proportion of different features in training set is the same as the proportion in testing set. And for each sampling rate, 10 different training set will be randomly selected by feeding various random seed, in order to demonstrate the impact of training set selection.}

In general, regression does not work well if the values fall into a wide range, since most regression
algorithms target minimizing the root-mean-square error (RMSE) between predicted and real values, and
RMSE is biased toward high values (e.g., 10\% error at a high value affects RMSE more than
10\% error at a low value). To address this problem, instead of trying to predict the raw throughput,
we try to predict the throughput per server core, which has a smaller value range.

We first try linear regression, and find it does not work very well in general, with
a coefficient of determination ($R^2$) value lower than 0.6 in most cases. This is not surprising since many systems
have non-linear throughput with regard to certain parameters (Section~\ref{sec:cases}).

We then try a classic neural network based multilayer perceptron model (MLP)~\cite{MLP1998}. 
In our experiments, we use two hidden layers with 10 and 5 neural nodes, respectively. 
The weights between each layer are updated during the backpropagation by
an optimization algorithm. We use the Limited-memory Broyden–Fletcher–Goldfarb–Shanno
(L-BFGS) algorithm~\cite{L-BFGS1989} as our optimizer, which performs better for small datasets.
Our models are trained with 10,000 iterations. We use this setting because it 
performs better than other settings we have tried, though note that this work
is by no means an exhaustive search of the best model or setting.

%\todo{Yujie, provide
%some details of MLP.} \todo{An MLP model consists one input layer, one output layer, and multiple hidden layers. Each layer is fully connected with %its input and output layer. The number of nodes of input layer is determined by the number of experimental parameters. }
For non-linear models, we meet a challenge that there is no standard metric, like $R^2$ for linear
models, to tell whether the prediction result is good or not: multiple articles argue $R^2$ is inappropriate
for non-linear models~\cite{R2NonLinear,KvaLsethR2}; RMSE value depends on the value range and thus it is impossible to
set a threshold. We finally decide to use $\frac{RMSE}{MAX\_VALUE-MIN\_VALUE}$ 
% (typically max-min but min is usually much smaller than max in
% our experiments) 
as a
normalized RMSE value~\cite{normalizeRMSE}. 





\begin{figure*}[ht]
     \centering

    
     \begin{subfigure}[t]{0.32\textwidth}
         \centering
         \includegraphics[width=\textwidth]{fig/thread_norm/rmse/cicada_ycsb.pdf}
         \caption{Cicada YCSB (5221).}
         \label{fig:cicada-ycsb-r2}
     \end{subfigure}
    ~
     \begin{subfigure}[t]{0.32\textwidth}
         \centering
         \includegraphics[width=\textwidth]{fig/thread_norm/rmse/gam}
         \caption{GAM TPC-C (510).}
         \label{fig:gam-r2}
     \end{subfigure}
     ~
    \begin{subfigure}[t]{0.32\textwidth}
         \centering
         \includegraphics[width=\textwidth]{fig/thread_norm/rmse/herd}
         \caption{HERD YCSB (600).}
         \label{fig:herd-r2}
     \end{subfigure}
     
     %\begin{subfigure}[t]{0.32\textwidth}
     %    \centering
     %    \includegraphics[width=\textwidth]{fig/thread_norm/rmse/silo_tpcc_neworder}
     %    \caption{Silo TPC-C NewOrder (660)}
     %    \label{fig:silo-neworder-r2}
     %\end{subfigure}
     %~
    %\begin{subfigure}[t]{0.32\textwidth}
    %     \centering
    %     \includegraphics[width=\textwidth]{fig/thread_norm/rmse/star_ycsb}
    %     \caption{Star YCSB (2686)}
    %     \label{fig:star-ycsb-r2}
    % \end{subfigure}
    % ~
    % \begin{subfigure}[t]{0.32\textwidth}
    %     \centering
    %     \includegraphics[width=\textwidth]{fig/thread_norm/rmse/tapir.pdf}
    %     \caption{TAPIR YCSB (600)}
    %     \label{fig:tapir-r2}
    % \end{subfigure}
     
     \begin{subfigure}[t]{0.32\textwidth}
         \centering
         \includegraphics[width=\textwidth]{fig/thread_norm/rmse/drtm.pdf}
         \caption{DrTM TPC-C (342).}
         \label{fig:drtm-r2}
     \end{subfigure}
     ~
     \begin{subfigure}[t]{0.32\textwidth}
         \centering
         \includegraphics[width=\textwidth]{fig/thread_norm/rmse/janus}
         \caption{Janus TPC-C (1102).}
         \label{fig:janus-r2}
     \end{subfigure}
    ~
      \begin{subfigure}[t]{0.32\textwidth}
         \centering
         \includegraphics[width=\textwidth]{fig/thread_norm/rmse/silo_ycsb.pdf}
         \caption{Silo YCSB (144).}
         \label{fig:silo-ycsb-r2}
     \end{subfigure}
     
%   \begin{subfigure}[t]{0.32\textwidth}
%         \centering
%         \includegraphics[width=\textwidth]{fig/thread_norm/rmse/star_tpcc}
%         \caption{Star TPC-C (275)}
%         \label{fig:star-tpcc-r2}
%     \end{subfigure}
     
     \caption{Examples of systems with accurate prediction. For each sampling rate, we try 10 times with different training sets, each corresponding to
     a point in these figures. The number
     in parenthesis is the total number of data points we have collected. Silo TPC-C, Star TPC-C, Star YCSB, and TAPIR YCSB figures have
     similar trends and can be found at our anonymous website~\cite{FigureSite}.}
     \vspace{-.1in}
     \label{fig:prediction_good}
\end{figure*}

%begin{figure*}[ht]

     
%     \caption{Systems with OK prediction accuracy. For each sampling rate, we try 10 times with different training sets \todo{Yujie: add the total number of data points for each figure}.}
%     \label{fig:prediction-ok}
%\end{figure*}

\begin{figure*}[ht]
     \begin{subfigure}[t]{0.32\textwidth}
         \centering
         \includegraphics[width=\textwidth]{fig/thread_norm/rmse/cicada_tpcc.pdf}
         \caption{Cicada TPC-C (106).}
         \label{fig:cicada-tpcc-r2}
     \end{subfigure}
     ~
     \begin{subfigure}[t]{0.32\textwidth}
         \centering
         \includegraphics[width=\textwidth]{fig/thread_norm/rmse/calvin.pdf}         
         \caption{Calvin (499).}
         \label{fig:calvin-r2}
     \end{subfigure}
~
    \begin{subfigure}[t]{0.32\textwidth}
         \centering
         \includegraphics[width=\textwidth]{fig/thread_norm/rmse/aria.pdf}
         \caption{Aria (605).}
         \label{fig:aria-r2}
     \end{subfigure}
     
     \caption{Systems with bad prediction accuracy. For each sampling rate, we try 10 times with different training sets, each corresponding to
     a point in these figures. The number
     in parenthesis is the total number of data points we collected.}
     \vspace{-.1in}
     \label{fig:prediction-bad}
\end{figure*}


%Figure~\ref{fig:prediction_good} and Figure~\ref{fig:prediction-bad} show the results.
Figure~\ref{fig:prediction_good} shows that most systems can achieve good
accuracy (normalized RMSE smaller than or close to 0.1) given enough samples,
which indicates sampling and regression is a promising idea to reduce experiment
cost.

Figure~\ref{fig:prediction-bad} shows the three systems+benchmarks that have 
a higher normalized RMSE and we investigate the reasons. For Cicada TPC-C,
the reason is that there are too few total data points (i.e., 106) as we can see the normalized
RMSE of the prediction is still declining with more samples. In practice, this should not be a problem
since we can simply run all experiments in this case. For Calvin, our investigation
shows that it has performance anomalies at a number of settings due to implementation
issues (Section~\ref{sec:cases}).
%\todo{Since Aria's authors implement a new version of
%Calvin, we try the new version as well and find ......}
Aria is the most interesting case because, first, as shown
in Figure~\ref{fig:aria-r2}, some training sets can achieve very good accuracy but
some do not, and second, our manual investigation finds its performance pattern is
very regular. We tried different models and feature settings, and found
that removing \#threads from the feature list can significantly improve the prediction
accuracy: in Figure~\ref{fig:aria-r2}, with a 50\% sampling rate, five training sets lead
to normalized RMSE over 0.2 but after removing the \#threads feature, only one training set
leads to normalized RMSE over 0.2. Therefore we suspect its inaccuracy is caused by sub-optimal model and feature
selection, which may cause problems like over-fitting. This may be addressed by
running multiple models together and using the best one.


%despite using a small number of samples; on GAM, Janus, Silo, Star TPC-C, we can get good accuracy with sufficient
%number of samples; on Aria, Calvin, Cicada TPC-C, the prediction is less inaccurate. 

%We further investigate two correlated questions about inaccurate prediction:
%1) for some systems, different sample sets with the same sampling rate can
%have a drastic difference in terms of accuracy. What is the reason behind that? 2) why are some systems
%less predictable significantly?

%Our investigation reveals two reasons for inaccurate prediction: first, if 
%we have too few total data points (e.g., Cicada TPC-C), then of course prediction
%is unlikely to be accurate. This should not be a problem in practice since we can
%simply run all experiments in this case. Second, we find some systems have
%abnormal performance under certain settings, which can explain both 1) and 2):
%if the training sets happen to include the abnormal points, the overall
%prediction results tend to be good; if the system has a lot of abnormal points,
%then it is less predictable than others.

%We will discuss the details of these abnormal points in Section~\ref{}, and
%we show some examples here.

%Calvin example: in this example, one can see Calvin with setting XX can have
%a very low throughput, while the settings around it has ``normal'' throughput.
%Show the prediction result after manually removing outliers.

%Maybe another example: ......

%Using MLP to automatically identify these points?

 

\vspace{-.1in}
\paragraph{Suggestions.}
In summary, sampling and regression could reduce total experiments
to run if the system has predictable performance; it could help to identify performance anomalies
for less predictable systems.

Based on our study, we
make the following suggestion: we can gradually test more settings and use maybe half of the results
to predict the other half, until the prediction accuracy is good enough
or becomes stable.
In these cases, we can be confident that we have tested enough settings to capture
the performance characteristics of the whole system.
From our study, we suggest to use normalized RMSE as the metric and use
0.1 as a threshold, though we acknowledge both
are debatable.
If the prediction accuracy stays low at the end, it is usually an indicator of performance
anomalies or inappropriate models, which requires further investigation.

\vspace{-.1in}
\paragraph{Remaining Questions.}
This study does not answer what parameters to tune and what values to choose,
which probably require a community-wide discussion. However, as a common software testing principle,
testing the extreme
values is usually a good idea to expose potential issues, as shown in Section~\ref{sec:cases}.

This study also does not address the problem that we may not have enough machines.
There are multiple works trying to evaluate a large cluster on a
limited number of machines~\cite{Gupta08DieCast,Agrawal11Emulating,Wang14Exalt,Shi2018Extrapolation,Stuardo2019ScaleCheck}, which may be helpful to some extent. 



%Note that this work is by no means an exhaustive search of or a proposal of
%the best regression model. Instead, the goal of this work is to show that, given
%we can only run a limited number of experiments, we can determine whether running more will
%have a significant marginal benefit by using regression: if the RMSE value does not drop
%any more, then probably we have reached our limit.